\documentclass[12pt]{article}

\usepackage{xelatex-generic,
			xelatex-biblatex,
			xelatex-glossaries
			}
\usepackage{soul,
			xcolor}
\sethlcolor{yellow}

\setotherlanguage{telugu}
\newfontfamily\telugufont
[Script=Telugu]
{Mandali}
			
\addbibresource{ss-editions.bib}

\title{Editions of the \emph{Suśrutasaṃhitā}}
\author{Deepro Chakraborty}
\date{\today}

\begin{document}
\maketitle



\begin{enumerate}
	\item \textbf{Tyler, John (1834) “Translation of two chapters of the First Part of the Soosroota”}.\\
	%
	The earliest known translation of any portion of the \emph{Suśrutasaṃhitā} into a European language appears to be John Tytler's “Translation of two chapters of the First Part of the Soosroota.” The work consists of a translation of chapters 4 and 5 of the \emph{Sūtrasthāna}. Tytler enclosed the translation in a letter dated 29 January 1834 addressed to Captain Troyer, secretary of the Sanskrit College from 1833 to 1835. The translation is reproduced in \fullcite[ch.,4][160--161]{Sen1991}.
	
	The translation was produced in the context of debates over medical education in colonial India. Following a proposal submitted by the Sanskrit College, Calcutta, to the General Committee of Public Instruction in June 1826, medical instruction was introduced at the College under Tytler's supervision. Students were taught both traditional Āyurvedic works, including the compendia of Caraka and Suśruta, and European medical texts. In March 1834, Captain Troyer sought Tytler's opinion on a proposal to abandon Sanskrit medical texts altogether and teach European medicine exclusively through English textbooks. By that time, Tytler had been serving as superintendent of the School for Native Doctors since 1830. As a committed orientalist and admirer of Sanskrit and Arabic scientific literature, he strongly opposed the proposal. The letter containing his translation of \emph{SS} 1.4--5 was written only two months before this exchange.\footnote{See \cite[145, 151--152, 160--161]{Sen1991}.}
	
	
	
	
	\item \textbf{\fullcite{Gupta1835-1836}}.\\ 
	The editio princeps of the \emph{Suśrutasaṃhitā} was published by Madhusudan Gupta, professor of medicine at the Sanskrit College, Calcutta. Gupta had also assisted John Tytler at the School for Native Doctors for two years, helping explain English medical terminology to students in Sanskrit and Bengali. The edition was published in two volumes in Devanagari script in 1835 and 1836.
	
	Gupta did not indicate the manuscripts or other sources used in preparing the text. The edition was reprinted in four volumes in Calcutta in 1868, and a second edition in two volumes appeared in 1874.\fvolcite{1B}[311]{HIMT}
	
	
	\item \textbf{\fullcite{Hessler1844}}.\\
	The German physician and philologist Franz Hessler translated the \emph{Suśrutasaṃhitā} into Latin and published the translation in three volumes in 1844, 1847, and 1850, respectively. He did not specify the edition or manuscripts on which his translation was based. 
	It is likely that he used Madhusudan Gupta's editio princeps as the basis of his translation, although he may also have consulted manuscripts. This possibility is suggested by his announcement, at the end of the preface to the first volume, that he intended to publish a “purified Sanskrit text” (\emph{textum Sanskrǐtum purgatum}) in two volumes.\fvolcite{1}[Praefatio VIII]{Hessler1844} This announcement suggests that he intended to establish a revised Sanskrit text, although it does not by itself demonstrate that he had consulted manuscript sources. No such edition, however, was ever published.
	In addition to the translation, Hessler later published two volumes of notes and annotations on the \emph{Suśrutasaṃhitā}.\footcite{Hessler1852}
	
	
	
	
	\item \textbf{\fullcite{Vidyasagara1873}}.\\
	Jīvānanda Vidyāsāgara Bhaṭṭācāryya published an edition of the \emph{Suśrutasaṃhitā} in Devanagari script at Calcutta in 1873. In many of his editions, Bhaṭṭācāryya reproduced texts from earlier printed editions without acknowledging his sources. 
	His edition of the \emph{SS} is also likely taken from Gupta's editio princeps.
	Subsequent editions were published in 1886, 1889, and 1899.
	
	
	
	
	\item \textbf{\fullcite{Bandopadhyaya1873}}.\\
	In Śaka 1795 (1873/74 \AD), Ambikā Caraṇa Bandyopādhyāya published an edition of the \emph{Suśrutasaṃhitā} in Bengali script together with a Bengali translation in two volumes. A second edition appeared in 1885. As in the case of Bhaṭṭācāryya's edition, no information is given regarding the sources used in establishing the text. It is likewise quite plausible that the edition ultimately derives from Gupta's editio princeps. 
	
	
	\item \textbf{\fullcite{Sen1880}}.\\
	\hl{[Not accessible to me.]}
	
	
	\item \textbf{\fullcite{Telugu1885}}.\\
	\hl{[Not yet accessible to me.]}
	
	
	\item \textbf{\fullcite{Kaviratna1885}}.\\
	\hl{[Not yet accessible to me.]}
	
	\item \textbf{\fullcite{Sena1886}}.\\
	\hl{[Not yet accessible to me.]} 
	\textcite[xxviii]{Cattopadhyaya1959} mentions that Niśikānta Sena could not complete this edition. 
	\citeauthor{Cattopadhyaya1959} tried to trace the 
	manuscript of Cakrapāṇidatta's \emph{Bhānumatī}, 
	which he believed was with someone in the Kumartuli area of Calcutta. But he was unsuccessful.
	
	
\end{enumerate}

\printbibliography

\end{document}
